%=======================================================================
%  COL 334 -- Assignment 2 : The Socket Exchange
%  Experiment Report
%
%  Compile with pdfLaTeX (Overleaf default).  Requires the folder
%  "figures/" to sit next to this file (see README-figures.txt).
%=======================================================================
\documentclass[11pt,a4paper]{article}

\usepackage[T1]{fontenc}
\usepackage[utf8]{inputenc}
\usepackage{lmodern}
\usepackage{amsmath}
\usepackage[margin=2.3cm]{geometry}
\usepackage{graphicx}
\usepackage{float}
\usepackage{booktabs}
\usepackage{longtable}
\usepackage{array}
\usepackage{xcolor}
\usepackage{listings}
\usepackage{caption}
\usepackage{enumitem}
\usepackage{microtype}
\usepackage{fancyhdr}
\usepackage[hidelinks]{hyperref}
\usepackage[most]{tcolorbox}   % tcolorbox asks to be loaded after hyperref

\graphicspath{{figures/}}

%----------------------------------------------------------------------
%  Styling
%----------------------------------------------------------------------
\definecolor{termbg}{RGB}{246,246,248}
\definecolor{termrule}{RGB}{200,200,210}
\definecolor{accent}{RGB}{25,70,130}

\lstdefinestyle{term}{
  basicstyle=\ttfamily\scriptsize,
  backgroundcolor=\color{termbg},
  frame=single,
  rulecolor=\color{termrule},
  breaklines=true,
  breakatwhitespace=false,
  columns=fullflexible,
  keepspaces=true,
  showstringspaces=false,
  xleftmargin=0.6em,
  xrightmargin=0.2em,
  framexleftmargin=0.4em,
  aboveskip=0.9em,
  belowskip=0.9em
}
\lstset{style=term}

\captionsetup{font=small,labelfont={bf,color=accent}}

\pagestyle{fancy}
\fancyhf{}
\fancyhead[L]{\small COL 334 -- Assignment 2: The Socket Exchange}
\fancyhead[R]{\small Experiment Report}
\fancyfoot[C]{\small\thepage}
\renewcommand{\headrulewidth}{0.4pt}

% Screenshot helper: #1 = file, #2 = caption, #3 = label
\newcommand{\shot}[3]{%
  \begin{figure}[H]
    \centering
    \fbox{\includegraphics[width=\linewidth,height=0.40\textheight,keepaspectratio]{#1}}
    \caption{#2}
    \label{#3}
  \end{figure}%
}

% Slightly shorter screenshot (for very wide, short captures)
\newcommand{\shotw}[3]{%
  \begin{figure}[H]
    \centering
    \fbox{\includegraphics[width=\linewidth,height=0.28\textheight,keepaspectratio]{#1}}
    \caption{#2}
    \label{#3}
  \end{figure}%
}

\newcommand{\qbox}[1]{%
  \begin{quote}\itshape #1\end{quote}%
}

% literal newline character, written as a C escape
\newcommand{\nlch}{\texttt{'\textbackslash n'}}

\newtcolorbox{answerbox}{
  colback=termbg, colframe=accent,
  boxrule=0.7pt, arc=2pt,
  left=7pt, right=7pt, top=6pt, bottom=6pt,
  breakable
}
\newenvironment{answer}
  {\begin{answerbox}\textbf{\color{accent}Answer.}\ }
  {\end{answerbox}}

%----------------------------------------------------------------------
\begin{document}

\begin{titlepage}
\centering
\vspace*{2.5cm}

{\Huge\bfseries The Socket Exchange\par}
\vspace{0.6cm}
{\LARGE Experiment Report\par}
\vspace{0.4cm}
{\large COL 334 -- Computer Networks \\ Assignment 2\par}

\vspace{2.2cm}

\begin{tabular}{@{}ll@{}}
\toprule
\textbf{Implementation language} & C (POSIX sockets, FreeBSD libc) \\
\textbf{Concurrency / I/O model} & Single-threaded, non-blocking, \texttt{kqueue}/\texttt{kevent} \\
\textbf{Platform}                & FreeBSD 14 (amd64) virtual machine, host \texttt{ragnarok} \\
\textbf{Transport}               & TCP over the loopback interface \texttt{lo0}, \texttt{127.0.0.1:5000} \\
\textbf{Harness}                 & \texttt{python3 experiment.py <n>} (provided with the assignment) \\
\bottomrule
\end{tabular}

\vspace{2.2cm}

\begin{minipage}{0.85\linewidth}
\small
This report documents Experiments~1--8 of Section~6 of the assignment handout.
For every experiment it states the question, the investigation approach, the exact
FreeBSD commands and tools used, the raw evidence (terminal screenshots and captured
output), and the answer supported by that evidence.  Section~\ref{sec:impl} describes
the concurrency and I/O design of the Exchange Server, which several of the
experiments probe directly.
\end{minipage}

\vfill
\end{titlepage}

\tableofcontents
\newpage

%======================================================================
\section{Experimental Setup}
%======================================================================

All processes -- the Exchange Server, the harness-generated clients and the
observation tools -- run inside a single FreeBSD virtual machine.  The server binds
\texttt{127.0.0.1:5000}, so every connection in this report is a loopback connection
and all packet captures are taken on \texttt{lo0}.

\paragraph{Launchers.}  The harness always starts the server through the standard
launcher, so it never needs to know the implementation language:

\begin{lstlisting}
$ cat server/run-server
#!/bin/sh
exec "$(dirname "$0")/../exchange_server" "$@"

$ cat client/run-trader
#!/bin/sh
exec "$(dirname "$0")/../trader_client" "$@"

$ cat client/run-market-data
#!/bin/sh
exec "$(dirname "$0")/../market_data_client" "$@"
\end{lstlisting}

\paragraph{Running an experiment.}  Each experiment is started from the submission
root with

\begin{lstlisting}
$ python3 experiment.py <experiment-number>
\end{lstlisting}

and observed from one or more additional SSH sessions into the same VM while the
harness is in its observation phase.

\paragraph{Notation in the command listings.}  Throughout this report a line beginning
with \texttt{\$} is a command that was run; a line beginning with \texttt{\#} and
followed by prose is a comment added here for the reader, not part of the command.  The
packet captures (\texttt{tcpdump}), the syscall traces (\texttt{ktrace}) and the kernel
tuning (\texttt{sysctl}) all require privilege and were run as \texttt{root}; the
screenshots therefore show a \texttt{\#} shell prompt for those sessions.

\paragraph{A note on the harness's readiness probe.}  Before every experiment the
harness opens a short-lived TCP connection to check that the server is listening and
then closes it abortively (\texttt{SO\_LINGER} with a zero timeout).  This probe is
visible in several packet captures as a connection that is answered with an
\texttt{RST}, and in the server's syscall traces as
\texttt{accept()} returning \texttt{errno 53 (ECONNABORTED)}.  It is not part of any
experiment; it is pointed out where it appears so that it is not mistaken for
experimental traffic.

%======================================================================
\section{Implementation Decisions}
\label{sec:impl}
%======================================================================

\subsection{Concurrency and I/O approach}

The Exchange Server is a \textbf{single process with a single thread} that runs one
\textbf{event loop built on FreeBSD's \texttt{kqueue}/\texttt{kevent}} interface.  There
is no thread and no process per connection.  Concretely:

\begin{itemize}[leftmargin=1.4em,itemsep=2pt]
  \item The listening socket (\texttt{fd 3}) is created with \texttt{socket()},
        \texttt{setsockopt(SO\_REUSEADDR)}, \texttt{bind()} and \texttt{listen()}, and is
        registered on the kqueue (\texttt{fd 4}) with an \texttt{EVFILT\_READ} filter.
  \item \textbf{Every} socket -- listening and accepted -- is put into non-blocking mode
        with \texttt{fcntl(fd, F\_SETFL, O\_NONBLOCK)}.  The server therefore never
        blocks inside \texttt{accept()}, \texttt{read()} or \texttt{write()}; the only
        blocking call in the whole program is \texttt{kevent()}.
  \item Each accepted socket is registered with \emph{two} filters,
        \texttt{EVFILT\_READ} and \texttt{EVFILT\_WRITE}, both with the
        \texttt{EV\_CLEAR} flag, i.e.\ edge-triggered notification.
  \item Because the notification is edge-triggered, \texttt{accept()} is called in a
        loop until it returns \texttt{EAGAIN}, so a burst of simultaneous connections is
        drained in one pass.  \texttt{ECONNABORTED} is tolerated and simply skipped
        (this is what happens to the harness's abortive readiness probe).
\end{itemize}

A symbol scan of the server binary agrees with the traces: it imports
\texttt{kqueue} and \texttt{kevent} and contains \emph{no} reference to
\texttt{poll}, \texttt{select} or any \texttt{pthread} function.  The event-loop
structure is directly visible in the \texttt{ktrace}/\texttt{kdump} evidence used in
Experiments~4--8; for example the accept path in Experiment~5:

\begin{lstlisting}
17538 exchange_server CALL  kevent(0x4,0,0,0x820883df0,0x40,0)
17538 exchange_server STRU  struct kevent[] = { { ident=3, filter=EVFILT_READ, ..., data=0x5, ... } }
17538 exchange_server RET   kevent 1
17538 exchange_server CALL  accept(0x3,0x820883d8c,0x820883db4)
17538 exchange_server STRU  struct sockaddr { AF_INET, 127.0.0.1:64427 }
17538 exchange_server RET   accept 5
17538 exchange_server CALL  fcntl(0x5,F_SETFL,0x6<O_RDWR|O_NONBLOCK>)
17538 exchange_server CALL  kevent(0x4,0x820883cc0,0x2,0,0,0)
17538 exchange_server STRU  struct kevent[] = { { ident=5, filter=EVFILT_READ,  flags=0x21<EV_ADD|EV_CLEAR>, ... }
                                               { ident=5, filter=EVFILT_WRITE, flags=0x21<EV_ADD|EV_CLEAR>, ... } }
...
17538 exchange_server CALL  accept(0x3,0x820883d8c,0x820883db4)
17538 exchange_server RET   accept -1 errno 35 Resource temporarily unavailable
\end{lstlisting}

\subsection{Why this approach}

\begin{itemize}[leftmargin=1.4em,itemsep=2pt]
  \item \textbf{No client can stall another.}  With one thread per connection, or with
        blocking I/O on a single thread, a client that opens a connection and then says
        nothing would park the server inside \texttt{recv()}.  With
        \texttt{kevent()} the server only ever touches a descriptor that the kernel has
        already declared ready, so an idle or half-speaking client costs nothing.  This
        is exactly what Experiment~4 measures.
  \item \textbf{Cost per connection is a few hundred bytes, not a stack.}  A
        thread-per-connection design pays a full thread stack (typically hundreds of
        kilobytes of virtual memory) per client.  Here a connection is one entry in a
        \texttt{realloc()}-grown table plus its I/O buffers.  Experiment~7 shows the
        server's RSS staying at exactly 2180~KiB while 5000 market-data updates are
        pushed to four clients.
  \item \textbf{\texttt{kqueue} is the native, scalable mechanism on FreeBSD.}  Unlike
        \texttt{select()} and \texttt{poll()}, it does not rescan the whole descriptor
        set on every call, and unlike \texttt{select()} it has no \texttt{FD\_SETSIZE}
        limit -- which matters for the bonus scalability target.
\end{itemize}

\subsection{Other design decisions relevant to TCP behaviour}

\paragraph{Message framing.}  The protocol is line-oriented, and TCP is a byte stream,
so every connection owns a growable input buffer.  After each \texttt{read()} the buffer
is scanned for \nlch; each complete line found is dispatched, and any trailing
partial line is compacted to the front of the buffer and kept until more bytes arrive.
The server therefore makes no assumption about how the sender split its
\texttt{send()} calls: one message may arrive in four reads (Experiment~3) and several
messages may arrive in one read.

\paragraph{Output buffering and back-pressure.}  Writes are non-blocking.  If a
\texttt{write()} to a client is short or returns \texttt{EAGAIN}, the remainder is
appended to that connection's own output buffer and the \texttt{EVFILT\_WRITE} event --
which is registered for every connection at accept time -- drains it later.  A slow
reader can therefore only make \emph{its own} output buffer grow; the event loop keeps
servicing everybody else.  This is what Experiment~7 verifies.

\paragraph{Roles.}  A connection starts out with no role.  The first command it sends
decides it: \texttt{LOGIN} makes it a Trader Client, \texttt{SUBSCRIBE} makes it a
Market-Data Client.  Afterwards, commands that do not belong to the role are rejected
with \texttt{ERROR}, which is how the capability matrix of Section~2.8 of the handout is
enforced.

\paragraph{Disconnect handling.}  A peer that closes gives \texttt{EVFILT\_READ} with
\texttt{EV\_EOF} and \texttt{read()} returning $0$.  The server then deletes both
filters for that descriptor, releases the connection slot (cancelling nothing else) and
calls \texttt{close()}.  One client disappearing never touches another client's state.

\paragraph{Graceful shutdown.}  \texttt{SIGINT}/\texttt{SIGTERM} interrupt the blocking
\texttt{kevent()} with \texttt{EINTR}; the loop then closes every connection slot, the
kqueue and the listening socket before exiting, which is visible at the end of every
\texttt{kdump} trace in this report.

%======================================================================
\section{Experiment 1 -- Listening and Connected Sockets}
%======================================================================

\qbox{After the client has connected, identify the TCP sockets associated with the
Exchange Server. How does the listening socket differ from the socket representing the
connection to the client?}

\subsection*{Approach and tools}

The harness (\texttt{python3 experiment.py 1}) starts the server and leaves exactly one
established, idle client connection in place.  While it waits, the server's sockets are
enumerated from two independent angles: \texttt{sockstat}, which maps sockets to the
owning process and file descriptor, and \texttt{netstat}, which shows the TCP state
machine's view.

\begin{lstlisting}
$ sockstat -4 -l                                   # listening IPv4 sockets
$ sockstat -4 -c                                   # connected IPv4 sockets
$ sockstat -4 -l | grep -E 'USER|exchange'         # same, filtered to the server
$ sockstat -4 -c | grep -E 'USER|exchange|python'
$ netstat -an -p tcp | grep 5000
\end{lstlisting}

\subsection*{Observations}

\shot{exp1/sockstat_all.jpeg}{\texttt{sockstat -4 -l} and \texttt{sockstat -4 -c}:
the full listening and connected IPv4 socket tables while the experiment client is
connected. The Exchange Server is PID~1874.}{fig:e1a}

\shot{exp1/sockstat_netstat.jpeg}{The same information filtered to the server and its
client, followed by \texttt{netstat -an -p tcp | grep 5000}: two \texttt{ESTABLISHED}
rows (one per endpoint of the loopback connection) and one \texttt{LISTEN}
row.}{fig:e1b}

\begin{lstlisting}
--- from: sockstat -4 -l | grep -E 'USER|exchange' -----------------------------
USER   COMMAND     PID  FD PROTO LOCAL ADDRESS      FOREIGN ADDRESS
harsh  exchange_s 1874  3  tcp4  127.0.0.1:5000     *:*                  <- listening

--- from: sockstat -4 -c | grep -E 'USER|exchange|python' ---------------------
USER   COMMAND     PID  FD PROTO LOCAL ADDRESS      FOREIGN ADDRESS
harsh  exchange_s 1874  5  tcp4  127.0.0.1:5000     127.0.0.1:47532      <- connected
harsh  python3.12 1873  3  tcp4  127.0.0.1:47532    127.0.0.1:5000       <- client end

--- from: netstat -an -p tcp | grep 5000 --------------------------------------
tcp4  0  0  127.0.0.1.5000   127.0.0.1.47532   ESTABLISHED
tcp4  0  0  127.0.0.1.47532  127.0.0.1.5000    ESTABLISHED
tcp4  0  0  127.0.0.1.5000   *.*               LISTEN
\end{lstlisting}

The Exchange Server (PID~1874) owns exactly two TCP sockets:

\begin{itemize}[leftmargin=1.4em,itemsep=2pt]
  \item \textbf{file descriptor 3}, local address \texttt{127.0.0.1:5000}, foreign
        address \texttt{*:*}, state \texttt{LISTEN};
  \item \textbf{file descriptor 5}, local address \texttt{127.0.0.1:5000}, foreign
        address \texttt{127.0.0.1:47532}, state \texttt{ESTABLISHED}.
\end{itemize}

The client process (\texttt{python3.12}, PID~1873) holds the mirror image of the second
socket, which is why \texttt{netstat} lists the connection twice: on loopback both
endpoints live on the same machine.

The gap in the numbering -- the first client gets \texttt{fd~5}, not \texttt{fd~4} --
is not an accident: \texttt{fd~4} is the kqueue the server waits on, which is a
descriptor but not a socket and therefore never appears in \texttt{sockstat} or
\texttt{netstat}.  It is visible in the descriptor table captured in
Experiment~5 (Figure~\ref{fig:e5c}).

\begin{answer}
The server has one \emph{listening} socket (fd~3) and one \emph{connected} socket
(fd~5).  The listening socket is only a half association -- local endpoint
\texttt{127.0.0.1:5000} with no peer (\texttt{*:*}) and state \texttt{LISTEN} -- and it
carries no data; its only role is to be handed to \texttt{accept()}.  The connected
socket is identified by the full four-tuple
\texttt{127.0.0.1:5000}\,$\leftrightarrow$\,\texttt{127.0.0.1:47532}, is in state
\texttt{ESTABLISHED} and is the descriptor on which the server actually reads and
writes.  \texttt{accept()} creates a brand-new descriptor for every client while the
listening socket keeps listening on the same port.
\end{answer}

%======================================================================
\section{Experiment 2 -- Observing TCP Connection States}
%======================================================================

\qbox{How does the TCP state of the client--server connection change during its
lifetime, and what events cause the observed changes?}

\subsection*{Approach and tools}

The harness (\texttt{python3 experiment.py 2}) connects a client, holds it idle for
10~seconds (phase~1), closes it, and keeps the server alive for another 15~seconds
(phase~2).  Two observers run in parallel: repeated \texttt{netstat} snapshots for the
state column, taken by hand at each announced phase of the run, and \texttt{tcpdump} on
\texttt{lo0} for the packets that cause the transitions.  The harness prints
\texttt{Phase~1} and \texttt{Phase~2} as it reaches them, which is what the snapshots
were timed against; because the two phases last 10 and 15~seconds there is no need to
poll faster.

\begin{lstlisting}
# terminal 2 -- the same command, re-run at each stage of the run:
#   twice during phase 1 (client connected and idle),
#   then repeatedly after phase 2 closed the connection,
#   and once more after the server itself exited
$ netstat -an -p tcp | grep 5000

# terminal 3 -- the packets that drive the state machine
$ tcpdump -i lo0 -n -S tcp port 5000
\end{lstlisting}

\subsection*{Observations}

\shot{exp2/harness.jpeg}{Harness output. Server PID~3225; the experiment client
connects from \texttt{127.0.0.1:25275}, stays idle for phase~1, and is closed at the
start of phase~2.}{fig:e2a}

\shot{exp2/tcpdump.jpeg}{\texttt{tcpdump -i lo0 -n -S tcp port 5000}: the three-way
handshake at 19:04:49, ten seconds of complete silence, then the four-way teardown at
19:04:59.}{fig:e2b}

\shot{exp2/netstat_states.jpeg}{\texttt{netstat -an -p tcp | grep 5000}, re-run by hand
at each stage -- eight invocations in all. The first produces no output because it was
run before the harness had started the server; the next two are taken during phase~1 and
show both endpoints \texttt{ESTABLISHED} alongside the \texttt{LISTEN} socket; the
following four are taken after phase~2 closed the connection and show only
\texttt{LISTEN}; the last one, after the server exited, is empty again.}{fig:e2c}

The capture in Figure~\ref{fig:e2b} contains the whole life of the connection
(client port 25275):

\begin{lstlisting}
19:04:49.840627  127.0.0.1.25275 > 127.0.0.1.5000  Flags [S],  seq 407608313
19:04:49.840656  127.0.0.1.5000  > 127.0.0.1.25275 Flags [S.], seq 1534167980, ack 407608314
19:04:49.840681  127.0.0.1.25275 > 127.0.0.1.5000  Flags [.],  ack 1534167981
                 ---- 10 s of phase 1: not a single packet ----
19:04:59.857904  127.0.0.1.25275 > 127.0.0.1.5000  Flags [F.], seq 407608314
19:04:59.857971  127.0.0.1.5000  > 127.0.0.1.25275 Flags [.],  ack 407608315
19:04:59.858238  127.0.0.1.5000  > 127.0.0.1.25275 Flags [F.], seq 1534167981
19:04:59.858274  127.0.0.1.25275 > 127.0.0.1.5000  Flags [.],  ack 1534167982
\end{lstlisting}

Mapping packets to states:

\begin{center}
\small
\begin{tabular}{@{}lll@{}}
\toprule
\textbf{Event on the wire} & \textbf{Client state} & \textbf{Server state} \\
\midrule
before anything                & \texttt{CLOSED}      & \texttt{LISTEN} (fd 3) \\
\texttt{SYN} sent              & \texttt{SYN\_SENT}   & \texttt{SYN\_RCVD} \\
\texttt{SYN,ACK} + final \texttt{ACK} & \texttt{ESTABLISHED} & \texttt{ESTABLISHED} (new fd) \\
idle phase (no packets)        & \texttt{ESTABLISHED} & \texttt{ESTABLISHED} \\
client \texttt{close()} sends \texttt{FIN} & \texttt{FIN\_WAIT\_1} & \texttt{CLOSE\_WAIT} \\
server \texttt{ACK}s the \texttt{FIN}      & \texttt{FIN\_WAIT\_2} & \texttt{CLOSE\_WAIT} \\
server sees EOF, \texttt{close()} sends \texttt{FIN} & \texttt{TIME\_WAIT}$^\dagger$ & \texttt{LAST\_ACK} \\
client \texttt{ACK}s          & (entry removed)      & \texttt{CLOSED} \\
\bottomrule
\end{tabular}

\smallskip
{\footnotesize $^\dagger$ Skipped on this system -- see the second bullet below.}
\end{center}

Two points are worth noting about what \texttt{netstat} actually shows.

\begin{itemize}[leftmargin=1.4em,itemsep=2pt]
\item \textbf{The idle phase produces no packets at all.}  A TCP connection that is
      established but unused is pure kernel state; nothing is sent to keep it alive
      (keep-alives are off by default).  This is why an idle client is invisible on the
      wire yet still fully \texttt{ESTABLISHED} in \texttt{netstat}.
\item \textbf{No \texttt{TIME\_WAIT} row appears after the close.}  The connection
      simply disappears and only the \texttt{LISTEN} socket remains, even though the
      client performed the active close and should therefore have entered
      \texttt{TIME\_WAIT} for $2\,\text{MSL}$.  This is a loopback-specific FreeBSD
      behaviour: the sysctl \texttt{net.inet.tcp.nolocaltimewait} makes an endpoint
      whose peer address is local skip the \texttt{TIME\_WAIT} state entirely, and on
      the FreeBSD~14 series it still defaults to~1 (it was flipped to~0 and deprecated
      only in later FreeBSD, for RFC~9293 compliance).  The value on this VM can be
      confirmed with \texttt{sysctl net.inet.tcp.nolocaltimewait}.  Experiment~6
      shows the same connection reported as \texttt{CLOSED} rather than
      \texttt{TIME\_WAIT} while its descriptor is still open, which is the other half of
      the same behaviour.
\end{itemize}

The two exchanges at 19:04:49.730 and 19:04:49.839 in Figure~\ref{fig:e2b} are the
harness's readiness probes and not part of the experiment.  They fail in two different
ways, which is worth reading carefully: the first (port 34294) is answered by the
server's host with an \texttt{RST} because nothing was listening on port 5000 yet -- a
refused connection that never existed -- whereas the second (port 42356) completes the
handshake and is then torn down by an \texttt{RST} sent \emph{by the probe itself}, which
is the harness's deliberate \texttt{SO\_LINGER}-zero close.  Only the second one tells
the harness that the server is up, after which it starts the experiment proper.

\begin{answer}
The connection walks the standard TCP state machine, and every transition is caused by
a packet or a system call, never by time.  \texttt{connect()} on the client and the
handshake \texttt{SYN} $\rightarrow$ \texttt{SYN,ACK} $\rightarrow$ \texttt{ACK} take
both ends from \texttt{CLOSED}/\texttt{LISTEN} to \texttt{ESTABLISHED}; the connection
then stays \texttt{ESTABLISHED} through ten completely silent seconds; the client's
\texttt{close()} emits a \texttt{FIN}, putting the client in \texttt{FIN\_WAIT\_1} and
the server in \texttt{CLOSE\_WAIT}; the server's \texttt{read()} returns~0, it calls
\texttt{close()} and emits its own \texttt{FIN} (\texttt{LAST\_ACK}); the client's final
\texttt{ACK} closes both ends.  Only the \texttt{LISTEN} socket survives.
\end{answer}

%======================================================================
\section{Experiment 3 -- TCP as a Byte Stream}
%======================================================================

\qbox{Does the server receive the message as one complete unit, or in multiple pieces?
Explain what this demonstrates about TCP vs.\ application-level message boundaries.}

\subsection*{Approach and tools}

The harness sends the single protocol message
\texttt{"LOGIN experiment\_trader\textbackslash n"} as
four separate \texttt{send()} calls -- \texttt{"LOGIN~"}, \texttt{"experiment"},
\texttt{"\_trader"}, \texttt{"\textbackslash n"} -- with a 200~ms pause between them.  The traffic is
captured with payloads printed in ASCII so that the split is visible byte for byte:

\begin{lstlisting}
$ tcpdump -i lo0 -n -A tcp port 5000
\end{lstlisting}

\subsection*{Observations}

\shot{exp3/tcpdump_setup.jpeg}{Start of the capture: the two readiness probes -- port
63659 refused with an \texttt{RST} from the server's host because the server was not
listening yet, then port 57740 connecting successfully and aborting itself with an
\texttt{RST} -- and finally the experiment connection from \texttt{127.0.0.1:16192}
completing its handshake.}{fig:e3a}

\shot{exp3/tcpdump_pieces.jpeg}{The four fragments of one application message, each in
its own TCP segment. The first three are acknowledged by a pure \texttt{ACK}; the fourth
-- the one carrying the newline -- is acknowledged on the server's \texttt{OK} reply
itself (\texttt{Flags [P.], seq 1:4, ack 25}), which is the first byte the server sends
in the whole exchange.}{fig:e3b}

\begin{lstlisting}
03:45:41.029576  16192 > 5000  Flags [P.], seq  1:7,  length  6   ...LOGIN
03:45:41.070223  5000  > 16192 Flags [.],  ack 7,     length  0
03:45:41.243816  16192 > 5000  Flags [P.], seq  7:17, length 10   ...experiment
03:45:41.278234  5000  > 16192 Flags [.],  ack 17,    length  0
03:45:41.456138  16192 > 5000  Flags [P.], seq 17:24, length  7   ..._trader
03:45:41.490000  5000  > 16192 Flags [.],  ack 24,    length  0
03:45:41.671446  16192 > 5000  Flags [P.], seq 24:25, length  1   ...(newline)
03:45:41.672246  5000  > 16192 Flags [P.], seq  1:4,  length  3   ...OK
03:45:41.710056  16192 > 5000  Flags [.],  ack 4,     length  0
\end{lstlisting}

The evidence has two independent parts:

\begin{enumerate}[leftmargin=1.6em,itemsep=2pt]
\item \textbf{The message crosses the network in four pieces.}  Four data segments of
      6, 10, 7 and 1 bytes carry one 24-byte application message.  Their sequence
      numbers are contiguous (\texttt{1:7}, \texttt{7:17}, \texttt{17:24},
      \texttt{24:25}) -- TCP is delivering a byte stream and has no idea that these 24
      bytes form one unit.  Each one is delivered to the server's socket as soon as it
      arrives, so the server's \texttt{read()} returns four times.
\item \textbf{The server nevertheless answers exactly once.}  \texttt{OK} is sent at
      \texttt{03:45:41.672246}, 0.8~ms after the fourth segment but roughly 640~ms after
      the first.  The server buffered the first three fragments, found no \nlch, and
      only dispatched \texttt{LOGIN experiment\_trader} when the newline finally
      arrived.  The acknowledgement pattern says the same thing from the other side:
      the first three segments draw a bare \texttt{ACK} and nothing else, while the
      fourth is acknowledged on the \texttt{OK} segment itself (\texttt{ack 25} riding
      on the 3-byte reply) -- the server had nothing to say until the delimiter
      completed the message.
\end{enumerate}

\begin{answer}
The server receives the message in \emph{multiple} pieces -- four segments of 6, 10, 7
and 1 bytes, the first three drawing a bare \texttt{ACK} and the fourth being
acknowledged on the reply itself -- and it still produces exactly one \texttt{OK}.  This is the essential point about TCP: it is a byte stream with no message
boundaries, so a sender's \texttt{send()} calls have no relationship to the receiver's
\texttt{recv()} calls.  Application-level framing must be reconstructed by the receiver,
which our server does by buffering bytes per connection and splitting them at the
\nlch\ delimiter.
\end{answer}

%======================================================================
\section{Experiment 4 -- One Client Should Not Stall the Others}
%======================================================================

\qbox{When Client~1 is idle, can the server still accept and service Client~2? Identify
the server operation determining the answer.}

\subsection*{Approach and tools}

The harness makes Client~1 send a \emph{deliberately incomplete} message --
\texttt{"LOGIN blocked\_client"} with no trailing newline -- and then go permanently
silent.  This is the harsher version of an idle client: the server has already received
bytes for a message that will never be completed.  Two seconds later Client~2 connects
and sends a complete message, and the harness times how long the reply takes.

The decisive question is \emph{which syscall the server is sitting in}, so the process
itself is inspected rather than only the network:

\begin{lstlisting}
$ sockstat -4 -c | grep -E 'USER|exchange'
$ netstat -an -p tcp | grep 5000
$ procstat -k $(pgrep -f exchange_server)         # kernel stack: what is it blocked on?
$ ps -H -o pid,tid,%cpu,wchan,comm -p $(pgrep -f exchange_server)
\end{lstlisting}

\subsection*{Observations}

\shot{exp4/harness.jpeg}{Harness output. Client~1 (port 61838) sends a partial message
and falls silent; Client~2 (port 58885) then logs in and is answered
\texttt{'OK'} in \textbf{0.005 seconds}. The server's own log line
\texttt{6: LOGIN active\_client} names the descriptor it served, and no such line ever
appears for Client~1.}{fig:e4a}

\shot{exp4/procstat_netstat.jpeg}{System-level evidence: both connections are
\texttt{ESTABLISHED} and owned by the server (fds 5 and 6); \texttt{procstat -k} shows
the single server thread parked in \texttt{kqueue\_scan}/\texttt{kqueue\_kevent}; and
\texttt{ps -H} reports one thread with \texttt{WCHAN = kqread}.}{fig:e4b}

\begin{lstlisting}
USER   COMMAND     PID  FD PROTO LOCAL ADDRESS   FOREIGN ADDRESS
harsh  exchange_s 1925  5  tcp4  127.0.0.1:5000  127.0.0.1:61838     <- silent client 1
harsh  exchange_s 1925  6  tcp4  127.0.0.1:5000  127.0.0.1:58885     <- client 2

                      Recv-Q Send-Q
tcp4 0 0 127.0.0.1.5000  127.0.0.1.58885 ESTABLISHED
tcp4 0 0 127.0.0.1.58885 127.0.0.1.5000  ESTABLISHED
tcp4 0 0 127.0.0.1.5000  127.0.0.1.61838 ESTABLISHED
tcp4 0 0 127.0.0.1.61838 127.0.0.1.5000  ESTABLISHED
tcp4 0 0 127.0.0.1.5000  *.*             LISTEN

  PID    TID COMM            TDNAME  KSTACK
 1925 100328 exchange_server -       mi_switch sleepq_catch_signals sleepq_wait_sig
                                     _sleep kqueue_scan kqueue_kevent kern_kevent_fp
                                     kern_kevent_generic sys_kevent amd64_syscall ...

  PID    LWP %CPU WCHAN  COMMAND
 1925 100328  0.0 kqread exchange_server
\end{lstlisting}

Three facts fall out of this:

\begin{itemize}[leftmargin=1.4em,itemsep=2pt]
  \item Client~2's round trip took \textbf{5~milliseconds} -- there is no stall of any
        kind.  Both clients are simultaneously \texttt{ESTABLISHED} and both are owned
        by the same single-threaded server process (fds 5 and 6).
  \item The server has \textbf{exactly one thread} (\texttt{TID 100328}), so it cannot
        be hiding the problem by dedicating a thread to the silent client.
  \item That one thread is blocked in \texttt{kqueue\_scan} with
        \texttt{WCHAN = kqread}, i.e.\ inside \texttt{kevent()} -- \emph{not} inside
        \texttt{recv()} on fd~5.  Client~1's incomplete message is sitting in the
        server's per-connection input buffer waiting for a newline that will never come,
        and that costs nothing because the server never asks the kernel for more bytes
        from fd~5 until the kernel says fd~5 is readable again.
\end{itemize}

Two smaller details in the same figures confirm where Client~1's bytes ended up.  First,
every \texttt{Recv-Q} in the \texttt{netstat} output is \texttt{0}, including fd~5's:
the kernel is holding nothing back, so the server has already \emph{read} the partial
\texttt{"LOGIN blocked\_client"} and is holding it in its own input buffer -- the wait is
happening in user space, on a delimiter, not in a syscall.  Second, the server's log in
Figure~\ref{fig:e4a} prints \texttt{6: LOGIN active\_client} for Client~2 and no
corresponding line for fd~5, which is the same fact seen from the application side: the
partial line was buffered but never dispatched as a command.

\begin{answer}
Yes -- Client~2 is accepted and answered \texttt{OK} in 0.005~s while Client~1 sits
silently mid-message.  The operation that decides this is the \textbf{blocking call at
the top of the event loop: \texttt{kevent()} on the kqueue}, not a per-connection
\texttt{recv()}.  \texttt{procstat -k} and \texttt{ps} both show the server's only
thread sleeping in \texttt{kqueue\_scan}/\texttt{kqread}, so it wakes up for whichever
descriptor becomes ready and reads only descriptors the kernel has already reported as
readable.  Had the server called a blocking \texttt{recv()} on Client~1's socket, its
kernel stack would instead show \texttt{soreceive}, and Client~2's \texttt{LOGIN} would
have gone unanswered until the harness gave up -- it waits 5~seconds for the reply and
would have reported \texttt{None} instead of \texttt{'OK'}.
\end{answer}

%======================================================================
\section{Experiment 5 -- Multiple Clients and I/O Multiplexing}
%======================================================================

\qbox{At a particular point, which client connections are ready for the server to
service, and what evidence supports this?}

\subsection*{Approach and tools}

The harness opens five simultaneous connections and then has only clients~1, 3 and~5
send an application message; clients~2 and~4 stay idle.  Readiness is a property of
the kernel's view of each socket, so it is captured from two directions:

\begin{lstlisting}
# 1. record every syscall the server makes, including the kevent return values
$ rm -f /tmp/exp5.ktr
$ while :; do pid=$(pgrep -n exchange_server); if [ -n "$pid" ]; then \
      ktrace -f /tmp/exp5.ktr -t cit -p $pid; echo "attached to $pid"; break; fi; done
$ ktrace -C                                     # stop tracing when done

# 2. read the trace back
$ kdump -f /tmp/exp5.ktr | grep -B2 -A6 kevent | less
$ kdump -f /tmp/exp5.ktr | grep -B6 'LOGIN client'
$ kdump -f /tmp/exp5.ktr | grep -E 'ident=[0-9]+, filter=EVFILT_READ, flags=0x20<EV_CLEAR>'

# 3. map the server's descriptors back to real sockets, and see queued bytes
$ PID=$(pgrep -n exchange_server); procstat -f $PID; ps -o pid,stat,wchan,comm -p $PID
$ netstat -an -p tcp | grep 5000

# 4. second run: freeze the server the instant it starts listening, so that the
#    bytes the clients send stay queued in the kernel and netstat can show them
$ while :; do pid=$(pgrep -n exchange_server); \
      if [ -n "$pid" ] && sockstat -4 -l | grep -q ':5000'; then \
          kill -STOP $pid; echo "stopped $pid"; break; fi; done
$ kill -CONT $(pgrep -n exchange_server)
\end{lstlisting}

\subsection*{Observations}

\shot{exp5/harness.jpeg}{Harness output for the traced run (server PID~17538). Clients
1, 3 and 5 connect from ports 64427, 43792 and 38472 and send \texttt{LOGIN}; clients 2
and 4 (ports 28067 and 65363) stay idle. The server logs each login against the
descriptor number: \texttt{5:}, \texttt{7:}, \texttt{9:}.}{fig:e5a}

\shot{exp5/ktrace_attach.jpeg}{Attaching \texttt{ktrace} to the server as soon as it
starts, and the beginning of the trace: five \texttt{accept()} calls returning fds 5--9,
each immediately set \texttt{O\_NONBLOCK} and registered on the kqueue with
\texttt{EVFILT\_READ} and \texttt{EVFILT\_WRITE}.}{fig:e5b}

\shot{exp5/procstat_fds.jpeg}{\texttt{procstat -f} for the server: fd~3 is the listening
socket \texttt{127.0.0.1:5000 *:0}, fd~4 is the kqueue, and fds~5--9 are the five client
connections in the same order the harness created them. \texttt{ps} shows the process
in \texttt{kqread}.}{fig:e5c}

\shot{exp5/kdump_ready.jpeg}{The decisive evidence: \texttt{kevent()} returns
\texttt{EVFILT\_READ} readiness for \texttt{ident=5}, \texttt{ident=7} and
\texttt{ident=9} only, each with \texttt{data=0xf} (15 bytes pending), and the server
reads exactly those descriptors. The final \texttt{grep} lists all three ready
descriptors together.}{fig:e5d}

The descriptor-to-client mapping from \texttt{procstat -f} (Figure~\ref{fig:e5c}) is
unambiguous:

\begin{center}
\small
\begin{tabular}{@{}llll@{}}
\toprule
\textbf{fd} & \textbf{Client} & \textbf{Peer port} & \textbf{Sent data?} \\
\midrule
3 & (listening socket) & \texttt{*:0}   & --- \\
4 & (kqueue)           & ---            & --- \\
5 & Client 1 & 64427 & yes \\
6 & Client 2 & 28067 & no (idle) \\
7 & Client 3 & 43792 & yes \\
8 & Client 4 & 65363 & no (idle) \\
9 & Client 5 & 38472 & yes \\
\bottomrule
\end{tabular}
\end{center}

and the kqueue reports readiness for precisely fds 5, 7 and 9:

\begin{lstlisting}
17538 exchange_server STRU  struct kevent[] = { { ident=5, filter=EVFILT_READ,  flags=0x20<EV_CLEAR>, data=0xf,    ... }
                                               { ident=5, filter=EVFILT_WRITE, flags=0x20<EV_CLEAR>, data=0xbf88, ... }
                                               { ident=6, filter=EVFILT_WRITE, flags=0x20<EV_CLEAR>, data=0xbf88, ... }
                                               { ident=7, filter=EVFILT_WRITE, flags=0x20<EV_CLEAR>, data=0xbf88, ... }
                                               { ident=8, filter=EVFILT_WRITE, flags=0x20<EV_CLEAR>, data=0xbf88, ... }
                                               { ident=9, filter=EVFILT_WRITE, flags=0x20<EV_CLEAR>, data=0xbf88, ... } }
17538 exchange_server RET   kevent 6
17538 exchange_server CALL  read(0x5,0x3bbbdd4180c1,0xa0)
17538 exchange_server GIO   fd 5 read 15 bytes
      "LOGIN client_1
      "
17538 exchange_server RET   read 15/0xf
17538 exchange_server CALL  write(0x1,0x3bbbdd409000,0x12)
17538 exchange_server GIO   fd 1 wrote 18 bytes
      "5: LOGIN client_1
      "
17538 exchange_server RET   write 18/0x12
17538 exchange_server CALL  write(0x1,0x3bbbdd409000,0x2d)
17538 exchange_server GIO   fd 1 wrote 45 bytes
      "User client_1 logged in from 127.0.0.1:64427
      "
17538 exchange_server RET   write 45/0x2d
17538 exchange_server CALL  write(0x5,0x20105b,0x3)
17538 exchange_server GIO   fd 5 wrote 3 bytes
      "OK
      "
17538 exchange_server RET   write 3
17538 exchange_server CALL  read(0x5,0x3bbbdd4180c1,0xa0)
17538 exchange_server RET   read -1 errno 35 Resource temporarily unavailable
...   (identically for ident=7 and ident=9; fds 6 and 8 are never read)
\end{lstlisting}

That first return needs reading carefully, because it carries six events, not three.
Exactly \emph{one} of them is an \texttt{EVFILT\_READ} event, on \texttt{ident=5}; the
other five are \texttt{EVFILT\_WRITE} events, one for every connection including the two
idle ones.  This is not a contradiction -- it is what write-readiness means.  A socket is
``ready to write'' whenever its send buffer has room, and all five send buffers are
empty, which is why each reports \texttt{data=0xbf88} (49\,032 bytes of free space).
Write-readiness therefore says nothing about whether a client has spoken; only
\texttt{EVFILT\_READ} does, and that filter fires for descriptors 5, 7 and~9 alone.  The
targeted \texttt{grep} in Figure~\ref{fig:e5d} makes the point without the noise: across
the whole trace, exactly three lines match
\texttt{filter=EVFILT\_READ, flags=0x20<EV\_CLEAR>}, and they are \texttt{ident=5},
\texttt{7} and~\texttt{9}.

\texttt{data=0xf} is the kernel telling the application how many bytes are waiting:
15, exactly the length of \texttt{"LOGIN client\_1\textbackslash n"}.  The two
intervening writes are to \texttt{fd~1}, the server's own standard output -- they are
the log lines visible in Figure~\ref{fig:e5a} and carry no network traffic; the only
socket write is the 3-byte \texttt{"OK\textbackslash n"} on \texttt{fd~5}.  Note also the
trailing \texttt{read()} returning \texttt{EAGAIN}: because the filters are
edge-triggered (\texttt{EV\_CLEAR}), the server keeps reading a ready descriptor until
the kernel says there is nothing left.

The same picture is visible from outside the process.  In a second run the server was
stopped with \texttt{kill -STOP} the moment it began listening, so it never consumed
anything: the bytes the three active clients sent stayed queued in the kernel receive
buffers, where \texttt{netstat} can show them.  (This is also why the harness afterwards
reports that the server had to be killed -- a stopped process cannot handle
\texttt{SIGTERM}.)  Worth being precise about what those sockets are: the server was
frozen before it ever called \texttt{accept()}, so the five connections exist only in
the kernel's accept queue and no server descriptor points at them yet.  The three-way
handshake and the receive buffering are done entirely by the TCP stack, which is exactly
why the byte counts are still visible even though the application is not running.

\shot{exp5/harness_stopped_run.jpeg}{The second run (server PID~9147): client~1 = 42196,
client~2 = 62300, client~3 = 13037, client~4 = 26978, client~5 = 53133.}{fig:e5e}

\shot{exp5/netstat_recvq.jpeg}{\texttt{netstat -an -p tcp}: the server-side rows for
ports 53133, 13037 and 42196 (clients 5, 3 and 1) carry \texttt{Recv-Q = 15}; the rows
for the idle clients 62300 and 26978 carry \texttt{Recv-Q = 0}.  The
\texttt{CLOSED} row is the leftover of the harness's readiness probe.}{fig:e5f}

\begin{answer}
At the moment the three active clients have spoken, exactly \textbf{three} of the five
connections are ready: the ones belonging to clients~1, 3 and~5.  Two independent pieces
of evidence agree.  (i)~\texttt{kevent()} returns \texttt{EVFILT\_READ} events only for
\texttt{ident=5}, \texttt{7} and~\texttt{9}, each carrying \texttt{data=0xf}
(15 bytes pending), and \texttt{procstat -f} maps those descriptors to peer ports 64427,
43792 and 38472 -- clients 1, 3 and 5.  (ii)~\texttt{netstat} shows \texttt{Recv-Q = 15}
on exactly those three server-side sockets and \texttt{Recv-Q = 0} on the two idle ones.
The idle connections~2 and~4 are perfectly healthy \texttt{ESTABLISHED} sockets; they are
simply never reported \emph{readable}, so the server never touches them and pays nothing
for them.  (They are reported \emph{writable}, because their send buffers are empty and
therefore always have room -- but write-readiness says nothing about whether a client has
sent anything.)
\end{answer}

%======================================================================
\section{Experiment 6 -- FIN vs.\ RST: Orderly and Abrupt Termination}
%======================================================================

\qbox{How does abrupt termination differ from orderly shutdown? Identify the TCP event
on the network and the corresponding socket behaviour.}

\subsection*{Approach and tools}

The harness runs two parts back to back on the same server.  In \textbf{Part~A} a client
connects and calls \texttt{shutdown(SHUT\_WR)}, the orderly close.  In \textbf{Part~B} a
second client sets \texttt{SO\_LINGER} with a zero timeout and calls \texttt{close()},
which makes the kernel abort the connection instead of closing it.  Three observers run
simultaneously:

\begin{lstlisting}
# packets
$ tcpdump -i lo0 -n -tttt -S tcp port 5000

# socket state, once per second with timestamps
$ while :; do date +%H:%M:%S; netstat -an -p tcp | grep 5000; echo ---; sleep 1; done

# what the server's syscalls see
$ rm -f /tmp/exp6.ktr
$ while :; do pid=$(pgrep -n exchange_server); if [ -n "$pid" ]; then \
      ktrace -f /tmp/exp6.ktr -t cit -p $pid; break; fi; done
$ kdump -f /tmp/exp6.ktr | grep -n -E 'EV_EOF|errno 54|ECONNRESET|CALL +close'
$ kdump -f /tmp/exp6.ktr | grep -A6 'EV_EOF'
$ kdump -f /tmp/exp6.ktr | grep -n -E 'RET +accept'
\end{lstlisting}

\subsection*{Observations}

\shot{exp6/tcpdump_fin_rst.jpeg}{\texttt{tcpdump} on \texttt{lo0} covering both parts.
Part~A (port 56436, 15:11:11) ends with a four-packet \texttt{FIN} exchange; Part~B
(port 38885, 15:11:21) ends with a single \texttt{RST} about 2~ms after its own
\texttt{SYN}.  The capture contains three \texttt{RST}s in total, and only the last one
is the \texttt{SO\_LINGER} reset under study: at 15:11:11.480553 the \emph{server} side
resets port 23897 because nothing is listening on 5000 yet, and at 15:11:11.590489 the
harness's second readiness probe (port 56301) completes a handshake and immediately
resets it.  Both precede Part~A and belong to the harness, not to the
experiment.}{fig:e6a}

\begin{lstlisting}
--- Part A: orderly close, client port 56436 -------------------------------
15:11:11.593527  56436 > 5000  Flags [S],  seq 751935905
15:11:11.593600  5000  > 56436  Flags [S.], seq 3975613210, ack 751935906
15:11:11.593676  56436 > 5000  Flags [.],  ack 3975613211
15:11:11.593872  56436 > 5000  Flags [F.], seq 751935906          <- shutdown(SHUT_WR)
15:11:11.593937  5000  > 56436  Flags [.],  ack 751935907          <- server ACKs the FIN
15:11:11.607184  5000  > 56436  Flags [F.], seq 3975613211         <- server close()
15:11:11.607269  56436 > 5000  Flags [.],  ack 3975613212

--- Part B: abortive close, client port 38885 ------------------------------
15:11:21.623193  38885 > 5000  Flags [S],  seq 2754814503
15:11:21.623253  5000  > 38885  Flags [S.], seq 243721720, ack 2754814504
15:11:21.623677  38885 > 5000  Flags [.],  ack 243721721
15:11:21.625160  38885 > 5000  Flags [R.], seq 2754814504, ack 243721721, win 0
\end{lstlisting}

\shot{exp6/netstat_after_fin.jpeg}{After the orderly close (Part~A): from 15:11:12 to
15:11:21 the connection remains listed once per second, in state \texttt{CLOSED},
because the client's descriptor is still open.}{fig:e6b}

\shot{exp6/netstat_after_rst.jpeg}{After the abortive close (Part~B): from 15:11:23
onwards only the \texttt{LISTEN} socket is left. The reset connection was destroyed
immediately and never appears in a single snapshot.}{fig:e6c}

\shot{exp6/kdump_eof.jpeg}{Server-side syscalls. The orderly close surfaces as
\texttt{EVFILT\_READ} with \texttt{EV\_EOF} and \texttt{read()} returning \texttt{0}
bytes; the abortive one never becomes a connection at all --
\texttt{accept()} fails with \texttt{errno 53, Software caused connection abort}. No
\texttt{ECONNRESET} appears anywhere in the trace.}{fig:e6d}

\begin{lstlisting}
--- orderly: the server is told politely that the peer is done -------------
11815 exchange_server STRU struct kevent[] = { { ident=5, filter=EVFILT_READ,
                                                 flags=0x8020<EV_CLEAR|EV_EOF>, fflags=0, data=0, ... } }
11815 exchange_server CALL read(0x5,0x735f92180c1,0xa0)
11815 exchange_server GIO  fd 5 read 0 bytes
      ""
11815 exchange_server RET  read 0
11815 exchange_server CALL kevent(0x4, ... EV_DELETE for ident=5 READ and WRITE ... )
11815 exchange_server GIO  fd 1 wrote 36 bytes
      "Client 127.0.0.1:56436 disconnected"
11815 exchange_server CALL close(0x5)

--- abortive: the reset arrives before the server ever owns a descriptor ---
11815 exchange_server RET  accept -1 errno 53 Software caused connection abort
11815 exchange_server RET  accept -1 errno 35 Resource temporarily unavailable
\end{lstlisting}

One field in that trace is worth singling out: the \texttt{EV\_EOF} event carries
\texttt{fflags=0}.  On FreeBSD, \texttt{kqueue} reports a socket error by putting the
errno into \texttt{fflags} alongside \texttt{EV\_EOF}, so \texttt{fflags=0} is the
kernel stating positively that this was a clean end of stream and not a reset.  A
connection torn down by an \texttt{RST} after it had been accepted would surface here as
\texttt{fflags=ECONNRESET}; none appears anywhere in the trace.

The contrast is complete:

\begin{center}
\small
\begin{tabular}{@{}p{0.24\linewidth}p{0.35\linewidth}p{0.35\linewidth}@{}}
\toprule
 & \textbf{Orderly (Part A)} & \textbf{Abrupt (Part B)} \\
\midrule
Trigger &
  \texttt{shutdown(SHUT\_WR)} / \texttt{close()} &
  \texttt{SO\_LINGER}\,=\,0 then \texttt{close()} \\
Packets &
  \texttt{FIN} $\rightarrow$ \texttt{ACK} $\rightarrow$ \texttt{FIN} $\rightarrow$ \texttt{ACK} (4) &
  a single \texttt{RST} (1) \\
Data already sent &
  delivered before the \texttt{FIN} is reported &
  discarded \\
Peer's socket &
  becomes readable, \texttt{read()} returns \texttt{0} &
  destroyed; a later \texttt{read}/\texttt{write} gives \texttt{ECONNRESET}/\texttt{EPIPE} \\
\texttt{netstat} &
  visible as \texttt{CLOSED} while the fd lives (Fig.~\ref{fig:e6b}) &
  gone instantly (Fig.~\ref{fig:e6c}) \\
Our server's reaction &
  \texttt{EV\_EOF}, deregister filters, log, \texttt{close()} &
  \texttt{accept()} returns \texttt{ECONNABORTED}, skipped \\
\bottomrule
\end{tabular}
\end{center}

\begin{answer}
An orderly shutdown is a negotiated, four-packet conversation
(\texttt{FIN}/\texttt{ACK} in each direction) that guarantees all previously sent bytes
are delivered; the peer's socket becomes readable and \texttt{read()} returns \texttt{0},
which is how our server learns of it (\texttt{EVFILT\_READ} with \texttt{EV\_EOF},
followed by \texttt{close()}).  An abrupt termination is a \textbf{single \texttt{RST}
segment}: the connection is torn down unilaterally with no acknowledgement, any queued
data is thrown away, and the entry vanishes from \texttt{netstat} immediately instead of
passing through the \texttt{FIN\_WAIT}/\texttt{CLOSED} sequence.  In this run the
\texttt{RST} arrived while the connection was still in the listen queue, so the server
saw it as \texttt{accept()} failing with \texttt{ECONNABORTED} (errno~53); had the
connection already been accepted, the next \texttt{write()} would have failed with
\texttt{EPIPE}/\texttt{ECONNRESET} instead.  Either way one client's abrupt death is a
per-connection error that the event loop absorbs without affecting anyone else.
\end{answer}

%======================================================================
\section{Experiment 7 -- Backpressure and the Slow Receiver}
%======================================================================

\qbox{How does the server's TCP connection to the slow client behave, and what evidence
shows whether this affects communication with other clients?}

\subsection*{Approach and tools}

The harness connects four clients -- a Market-Data Client that keeps draining its
socket, a Market-Data Client that \emph{never} reads, and two Traders -- and then drives
5000 matching \texttt{BUY}/\texttt{SELL} pairs, so the server must push 5000
\texttt{TRADE JNST 1 238} messages (17~bytes each) to \emph{both} subscribers.  To give
backpressure the best possible chance of appearing, the TCP socket buffers were first
shrunk to 8~KiB and auto-tuning was switched off.  The run was then instrumented with a
set of periodic samplers writing to \texttt{/tmp/exp7/}:

\begin{lstlisting}
# 00 -- shrink the socket buffers and switch auto-tuning off, to make the
#       connection reach its buffer limit as early as possible; then record them
$ sysctl net.inet.tcp.recvbuf_auto=0 net.inet.tcp.sendbuf_auto=0
$ sysctl net.inet.tcp.recvspace=8192 net.inet.tcp.sendspace=8192
$ sysctl net.inet.tcp.recvspace net.inet.tcp.sendspace \
         net.inet.tcp.recvbuf_auto net.inet.tcp.sendbuf_auto | tee /tmp/exp7/00_sysctl.txt

# 10 -- watch specifically for zero-window advertisements
$ tcpdump -i lo0 -n -tttt "tcp port 5000 and tcp[14:2] = 0" | tee /tmp/exp7/10_zerowindow.txt

# 20 -- per-connection queue depths, every 2 s
$ while :; do date +%H:%M:%S; netstat -an -p tcp | grep 5000; echo ---; sleep 2; done \
      | tee /tmp/exp7/20_netstat.txt

# 30 -- the run itself
$ python3 experiment.py 7 | tee /tmp/exp7/30_harness.txt

# 40 -- server memory / CPU / what it is blocked on, every 2 s
$ while :; do date +%H:%M:%S; ps -o pid,rss,vsz,%cpu,stat,wchan -p $PID; echo ---; sleep 2; done \
      | tee /tmp/exp7/40_server.txt

# protocol counters before and after, then diff
$ netstat -s -p tcp > /tmp/exp7/01_tcpstats_before.txt   # ... run ...
$ netstat -s -p tcp > /tmp/exp7/02_tcpstats_after.txt
$ diff /tmp/exp7/01_tcpstats_before.txt /tmp/exp7/02_tcpstats_after.txt \
      | tee /tmp/exp7/03_tcpstats_diff.txt
\end{lstlisting}

Every sampler is piped through \texttt{tee} rather than redirected, because the
screenshots below are the live terminal views of exactly these commands.  All of them
were run as \texttt{root} from a single shell and write into \texttt{/tmp/exp7/}, so the
files do not end up split across two home directories depending on who ran what.

\subsection*{Observations}

\shotw{exp7/harness.jpeg}{The four connections and the generated load: normal
Market-Data client on port 35877, \textbf{slow} Market-Data client on port 50266, buyer
on 65121, seller on 49148, and \textbf{5000 matching trades} generated.}{fig:e7a}

\shotw{exp7/sysctl.jpeg}{Socket-buffer parameters recorded at the start of the
run.}{fig:e7b}

\shot{exp7/netstat_queues.jpeg}{Queue depths at the peak. The slow client's own socket
(\texttt{127.0.0.1.50266}) holds \texttt{Recv-Q = 85000} bytes that its application has
never read, while the server's send queues to \emph{all} clients, including the slow
one, are essentially empty.}{fig:e7c}

\shotw{exp7/server_ps.jpeg}{The server (PID~19612) at the beginning and the end of the
run: RSS unchanged at 2180~KiB, CPU falling back to 0.0\,\%, and \texttt{WCHAN} still
\texttt{kqread}.}{fig:e7d}

\shotw{exp7/zerowindow.jpeg}{Searching the capture for zero-window advertisements. The
only two matches are \texttt{RST} segments (which always carry \texttt{win 0}); no
genuine zero-window advertisement was ever sent.}{fig:e7e}

The time series extracted from \texttt{20\_netstat.txt} is the core result.  Writing
$Q_{\text{recv}}$ for the slow client's own receive queue and $Q_{\text{send}}$ for the
server's send queue towards it:

\begin{center}
\small
\begin{tabular}{@{}lrrrrrrrr@{}}
\toprule
Time (\texttt{hh:mm:ss}) & 16:46:09 & 16:46:12 & 16:46:16 & 16:46:30 & 16:47:00 & 16:48:01 & 16:49:11 & 16:49:34 \\
\midrule
$Q_{\text{recv}}$ (slow client)  & 476 & 1\,343 & 3\,060 & 8\,959 & 21\,624 & 46\,750 & 76\,228 & 85\,000 \\
$Q_{\text{send}}$ (server side)  & 17  & 17    & 34    & 17    & 17     & 17     & 34     & 0 \\
\bottomrule
\end{tabular}
\end{center}

Over all 108 samples of the run:

\begin{itemize}[leftmargin=1.4em,itemsep=2pt]
\item The slow client's \texttt{Recv-Q} rises \textbf{monotonically} and levels off at
      \textbf{exactly 85\,000 bytes} $= 5000 \times 17$ -- every single \texttt{TRADE}
      message was delivered into its socket buffer and none of them was ever read by the
      application.\footnote{The plateau is an order of magnitude above the 8~KiB
      \texttt{net.inet.tcp.recvspace} configured for this run (Figure~\ref{fig:e7b}), so
      the kernel evidently sized this loopback socket's receive buffer well beyond the
      nominal value.  \texttt{recvspace} is only the \emph{initial} size; the hard
      ceiling is \texttt{kern.ipc.maxsockbuf}, which was left at its default and is
      megabytes wide, so 85~kB is comfortably inside what the socket is allowed to hold.
      We did not chase the exact sizing rule; what the measurement
      establishes either way is that the receiver's buffer was never exhausted, which is
      what the rest of the argument rests on.}
\item The server's \texttt{Send-Q} towards the slow client \textbf{never exceeded
      34~bytes} (two pending messages) across the entire run, and the same holds for the
      normal client.
\item \texttt{netstat -s -p tcp} confirms this from the protocol counters:
      \texttt{0 window probe packets}, \texttt{0 window probes} and
      \texttt{0 persist timeouts} -- the persist timer never ran, so no receiver ever
      advertised a closed window.
\item The server process was completely unaffected: RSS pinned at 2180~KiB for the whole
      run, CPU decaying from 1.1\,\% to 0.0\,\%, and \texttt{WCHAN} always
      \texttt{kqread} -- never \texttt{sbwait}, which is what a blocking
      \texttt{write()} would show.
\item The other three clients kept working throughout: all 5000 trades completed, the
      normal subscriber's queues never exceeded 34 bytes, and the two traders kept getting
      their \texttt{ORDER\_ACCEPTED}/\texttt{BOUGHT}/\texttt{SOLD} messages.
\end{itemize}

\paragraph{Interpretation.}  Even with the socket buffers deliberately reduced to 8~KiB,
the 85~KiB this experiment generates fit entirely inside the slow receiver's socket
buffer, so TCP's flow control never had to engage on loopback: the receiver kept
advertising a non-zero window, the server's writes all succeeded immediately, and the
backlog simply accumulated in the slow client's own receive buffer.  Had the stream been large enough to fill that buffer, the mechanism would have
been: receiver advertises \texttt{win 0} $\rightarrow$ the server's \texttt{Send-Q}
climbs to the send-buffer limit $\rightarrow$ \texttt{write()} returns \texttt{EAGAIN}
$\rightarrow$ the server appends to that connection's user-space output buffer and waits
for \texttt{EVFILT\_WRITE}.  Crucially, none of those steps involve blocking, so the
head-of-line blocking stays confined to the one slow connection.  A blocking-write
server would instead have parked in \texttt{write()} with \texttt{WCHAN = sbwait} and
stopped serving everyone.

\begin{answer}
The connection to the slow client stays \texttt{ESTABLISHED} and healthy; the unread
data piles up on the \emph{receiver's} side, its \texttt{Recv-Q} climbing monotonically
to 85\,000~bytes -- exactly the 5000 $\times$ 17~bytes of market data that the server
sent.  At this volume the receiver's socket buffer was never exhausted, so no zero
window was ever advertised (\texttt{0 window probes}, \texttt{0 persist timeouts} in
\texttt{netstat -s}), and the server's send queue to that client never exceeded
34~bytes.  Communication with the other clients was completely unaffected: all 5000
trades were generated, the normal subscriber's queues stayed near zero, and the server's
RSS (2180~KiB), CPU (0\,\%) and wait channel (\texttt{kqread}, i.e.\ idle in
\texttt{kevent()}) never changed.  Because every write is non-blocking and any residue
goes into a per-connection output buffer drained by \texttt{EVFILT\_WRITE}, a slow
reader can only ever penalise itself.
\end{answer}

%======================================================================
\section{Experiment 8 -- Unexpected Client Disconnection}
%======================================================================

\qbox{When a client disappears unexpectedly, what happens to the TCP connection, and
what network-level evidence lets you determine how the server detects the failure?}

\subsection*{Approach and tools}

The harness starts a \emph{separate process} that acts as a Market-Data Client,
subscribes to \texttt{JNST} and then does nothing but sleep, so the server is actively
sending it trade updates that it never reads.  After 20 trades the harness
\texttt{SIGKILL}s that process -- the application gets no chance to close anything.  A
second Market-Data Client and the two traders stay connected so that the two situations
can be compared side by side.  Packets, socket states and the server's syscalls were all
recorded to \texttt{/tmp/exp8/}:

\begin{lstlisting}
# 10 -- full packet capture of the port
$ tcpdump -i lo0 -n -tttt -S tcp port 5000            | tee /tmp/exp8/10_tcpdump.txt

# 20 -- socket table once per second, with timestamps
$ while :; do date +%H:%M:%S; netstat -an -p tcp | grep 5000; echo ---; sleep 1; done \
                                                      | tee /tmp/exp8/20_netstat.txt

# server syscalls -- armed before the run and left spinning, so that ktrace attaches
#   within milliseconds of the server appearing and the accept() calls are captured
$ while :; do pid=$(pgrep -n exchange_server); if [ -n "$pid" ]; then \
      ktrace -f /tmp/exp8/server.ktr -t cit -p $pid; echo "attached to $pid"; break; fi; done

# post-processing
$ grep -E 'Flags \[R|Flags \[F' /tmp/exp8/10_tcpdump.txt
$ kdump -f /tmp/exp8/server.ktr | grep -B1 -E 'RET +accept [0-9]'
$ kdump -f /tmp/exp8/server.ktr | sed -n '1620,1650p'
$ awk '/^17:10:02$/,/^17:10:06$/' /tmp/exp8/20_netstat.txt
\end{lstlisting}

\subsection*{Observations}

\shotw{exp8/harness.jpeg}{The three connections created directly by the harness:
surviving Market-Data client on port 10123, buyer on 38997, seller on 26991. The
client that will be killed is a separate process and connects on its own.}{fig:e8a}

\shotw{exp8/kdump_accept.jpeg}{\texttt{kdump} of the server's \texttt{accept()} calls
identifies the fourth connection: fd~5 = 10123 (surviving MD), fd~6 = 38997 (buyer),
fd~7 = 26991 (seller), \textbf{fd~8 = 48664 (the client that will disappear)}.}{fig:e8b}

\shot{exp8/netstat.jpeg}{\texttt{netstat} one second at a time. At 17:10:02 and 17:10:03
the doomed connection (\texttt{127.0.0.1.48664}) is \texttt{ESTABLISHED} with 136 and
then 255 bytes of unread market data in its \texttt{Recv-Q}. At 17:10:04 it is gone
completely, while the other three connections carry on exchanging data.}{fig:e8c}

\shot{exp8/tcpdump_fin.jpeg}{Every \texttt{FIN} and \texttt{RST} on the port. At
\texttt{17:10:03.772273} the killed client's kernel sends a \texttt{FIN}, the server
\texttt{ACK}s it and replies with its own \texttt{FIN} 4~ms later. The other three
connections only close at 17:10:23, when the experiment ends.}{fig:e8d}

\shot{exp8/kdump_eof.jpeg}{The server's view of the death: \texttt{kevent()} reports
\texttt{ident=8} with \texttt{EV\_CLEAR|EV\_EOF}, \texttt{read(8)} returns \textbf{0
bytes}, the filters are deleted, the disconnect is logged and \texttt{close(8)} is
called. The very next \texttt{kevent()} already returns readiness for \texttt{ident=6}
and \texttt{ident=7} -- the other clients were never interrupted.}{fig:e8e}

The complete sequence around the kill:

\begin{lstlisting}
--- packets ---------------------------------------------------------------
17:10:03.632762  5000  > 48664  Flags [P.], length 17     "TRADE JNST 1 238"
17:10:03.726335  48664 > 5000   Flags [.],  ack ...       (kernel ACKs; app never reads)
17:10:03.772273  48664 > 5000   Flags [F.], seq 584026206 <- process is SIGKILLed here
17:10:03.772319  5000  > 48664  Flags [.],  ack 584026207
17:10:03.776622  5000  > 48664  Flags [F.], seq 2245618379
17:10:03.776710  48664 > 5000   Flags [.],  ack 2245618380

--- socket table ----------------------------------------------------------
17:10:03  tcp4  255  0  127.0.0.1.48664  127.0.0.1.5000  ESTABLISHED
17:10:04  (no row for 48664 at all; 26991, 38997 and 10123 still ESTABLISHED
           with Send-Q 34, 36 and 17 bytes -- traffic continues normally)

--- server syscalls -------------------------------------------------------
31814 exchange_server STRU struct kevent[] = { { ident=8, filter=EVFILT_READ,
                                                 flags=0x8020<EV_CLEAR|EV_EOF>, fflags=0, data=0, ... } }
31814 exchange_server RET  kevent 1
31814 exchange_server CALL read(0x8,0x3172b7018541,0xa0)
31814 exchange_server GIO  fd 8 read 0 bytes
      ""
31814 exchange_server RET  read 0
31814 exchange_server CALL kevent(...)   EV_DELETE for ident=8 READ and WRITE
31814 exchange_server GIO  fd 1 wrote 36 bytes
      "Client 127.0.0.1:48664 disconnected"
31814 exchange_server CALL close(0x8)
31814 exchange_server CALL kevent(0x4,0,0,0x82085e0c0,0x40,0)
31814 exchange_server STRU struct kevent[] = { { ident=6, filter=EVFILT_READ, ..., data=0xf }
                                              { ident=7, filter=EVFILT_READ, ..., data=0x10 } }
\end{lstlisting}

Three things deserve emphasis.

\begin{itemize}[leftmargin=1.4em,itemsep=2pt]
\item \textbf{The connection was closed by the kernel, not the application.}  The client
      process was \texttt{SIGKILL}ed, so it executed no shutdown code at all.  FreeBSD
      still released its file descriptors on process exit, and the TCP stack performed a
      normal active close.  On the wire an abruptly killed \emph{process} therefore looks
      like an orderly close: a \texttt{FIN}, not an \texttt{RST}.
\item \textbf{Unread data is visible right up to the moment of death.}  \texttt{Recv-Q}
      on the doomed socket grows 17~$\rightarrow$~136~$\rightarrow$~255~bytes as trades
      arrive that the sleeping process never reads, and one second later the socket is
      gone entirely -- no \texttt{TIME\_WAIT} row is left behind, for the same
      \texttt{net.inet.tcp.nolocaltimewait} reason established in Experiment~2
      (page~\pageref{fig:e2b}): the peer address is local, so the endpoint that performed
      the active close skips \texttt{TIME\_WAIT} altogether.
\item \textbf{Detection is immediate and local.}  The server did not need a timeout, a
      keep-alive or a failed write: the \texttt{FIN} made the descriptor readable,
      \texttt{read()} returned~0, and the connection slot was released.  The
      \texttt{EV\_EOF} event again carries \texttt{fflags=0}, and searching the whole
      trace for \texttt{errno 32 (EPIPE)} or \texttt{errno 54 (ECONNRESET)} returns
      nothing -- the server never even attempted a write that could fail.  The next
      \texttt{kevent()} call in the same trace is already servicing fds~6 and~7, and the
      surviving subscriber received all 50 of the post-kill trades.
\end{itemize}

\paragraph{The harder failure mode.}  This experiment kills a process, and the OS is
still there to send the \texttt{FIN}.  If instead the \emph{machine} died or the network
were cut, no \texttt{FIN} and no \texttt{RST} would ever arrive; the server's socket
would stay \texttt{ESTABLISHED} indefinitely.  It would then only find out by writing --
retransmissions time out and the connection is dropped with \texttt{ETIMEDOUT}, or the
peer reboots and answers with an \texttt{RST}, giving \texttt{ECONNRESET}/\texttt{EPIPE}
-- or by enabling \texttt{SO\_KEEPALIVE} and letting
\texttt{net.inet.tcp.keepidle}/\texttt{keepintvl}/\texttt{keepcnt} probe the peer.

\begin{answer}
The client process vanished without running any shutdown code, but the FreeBSD kernel
closed its descriptors on exit, so the TCP connection was terminated with a normal
\textbf{\texttt{FIN} exchange} at \texttt{17:10:03.772}, not a reset -- the server
\texttt{ACK}ed and sent its own \texttt{FIN} 4~ms later, and the socket disappeared from
\texttt{netstat} within one second (having held 255~bytes of undelivered market data at
the moment of death).  The server detected the failure through that \texttt{FIN}: the
descriptor became readable with \texttt{EVFILT\_READ | EV\_EOF} and \texttt{read()}
returned \texttt{0} bytes, whereupon it deregistered the filters, logged
\texttt{Client 127.0.0.1:48664 disconnected} and closed fd~8.  The other three
connections were not disturbed at all -- the very next \texttt{kevent()} return in the
same trace services fds~6 and~7.
\end{answer}

%======================================================================
\section{Summary of Tools Used}
%======================================================================

\begin{center}
\small
\begin{tabular}{@{}llp{0.46\linewidth}@{}}
\toprule
\textbf{Tool} & \textbf{Experiments} & \textbf{What it contributed} \\
\midrule
\texttt{sockstat -4 -l / -c} & 1, 4 &
  Maps every socket to its owning process and file descriptor; separates listening from
  connected sockets. \\
\texttt{netstat -an -p tcp} & 1, 2, 4, 5, 6, 7, 8 &
  TCP state column, and the \texttt{Recv-Q}/\texttt{Send-Q} depths used as evidence that
  a partial message had already been consumed by the server (Exp.~4), of readiness
  (Exp.~5) and of backpressure (Exp.~7, 8). \\
\texttt{netstat -s -p tcp} & 7 &
  Protocol counters: \texttt{0 window probes} and \texttt{0 persist timeouts} proved no
  zero-window event occurred. \\
\texttt{tcpdump -i lo0} & 2, 3, 6, 7, 8 &
  The packet-level ground truth: handshake, per-segment framing (\texttt{-A}),
  \texttt{FIN} vs.\ \texttt{RST}, zero-window search. \\
\texttt{ktrace} + \texttt{kdump} & 5, 6, 8 &
  The server's own syscalls: \texttt{kevent()} return values naming the ready
  descriptors, \texttt{read()} returning 0 at EOF, \texttt{accept()} errno values. \\
\texttt{procstat -k} & 4 &
  Kernel stack of the server thread -- proved it sleeps in
  \texttt{kqueue\_scan}, not in \texttt{recv()}. \\
\texttt{procstat -f} & 5 &
  Descriptor table, used to map \texttt{kevent} idents back to real client sockets. \\
\texttt{ps -H -o \dots,wchan} & 4, 7 &
  Thread count, \texttt{WCHAN} (\texttt{kqread}), RSS and CPU of the server under load. \\
\texttt{sysctl} & 7 &
  Socket-buffer parameters that gate backpressure. \\
\bottomrule
\end{tabular}
\end{center}

%======================================================================
\appendix
\section{Bonus -- Connection Scalability and I/O Design}
\label{sec:bonus}
%======================================================================

This appendix reports the results of Section~6.9 of the handout: driving the Exchange
Server up to $N=70\,000$ simultaneous idle TCP connections, measuring the per-$N$
resource footprint, and analysing the bottlenecks reached.  All measurements were taken
on the same FreeBSD~14.4-RELEASE VM used for Experiments~1--8.  A companion document,
\texttt{bonus/README.md} in the submission, lists every OS-level tunable and command
verbatim; only the parts needed to read the tables and figures are repeated here.

%----------------------------------------------------------------------
\subsection{Setup and tooling}
%----------------------------------------------------------------------

Two small helper programs, both in \texttt{bonus/}, were written for this section.  The
Exchange Server, Trader Client and Market-Data Client themselves are otherwise unchanged
from Experiments~1--8 (with the single exception noted in \S\ref{sec:bonus-listen}).

\paragraph{\texttt{bonus/idle\_clients.py} -- client generator.}  Opens $N$ TCP
connections to the Exchange Server and holds them \texttt{ESTABLISHED} without ever
sending or receiving application data.  It raises \texttt{RLIMIT\_NOFILE} to $N+128$
before starting, reports progress and connect rate every \texttt{--batch} connections,
and records the index and \texttt{errno} of every \texttt{connect()} failure so the
count at which the server first fails is unambiguous.  Three additive optional flags
(\texttt{-{}-src-ips}, \texttt{-{}-src-cidr}, \texttt{-{}-src-count}) enable
source-IP rotation and are described in \S\ref{sec:bonus-srcip}; without them the
script behaves exactly like a default client (kernel-picked \texttt{127.0.0.1} +
ephemeral port).

\begin{lstlisting}
$ python3 bonus/idle_clients.py 127.0.0.1 5000 50000 --batch 2000
$ python3 bonus/idle_clients.py 127.0.0.1 5000 70000 --batch 2000 \
        --src-cidr 127.0.0.0/8 --src-count 4
\end{lstlisting}

\paragraph{\texttt{bonus/measure.sh} -- one snapshot per row of the table.}  Runs
while \texttt{idle\_clients.py} is holding its connections open and captures every
column of the required Section~6.9 table in a single output:

\begin{lstlisting}
$ ./bonus/measure.sh <server-pid> 5000 | tee run_<N>.txt

  established conns   netstat -an -p tcp | awk ... | wc -l  (+ sockstat cross-check)
  server FDs          procstat -f <pid> | tail -n +2 | wc -l
  memory / CPU        ps -o pid,rss,vsz,%cpu,%mem,nlwp,command -p <pid>
  system-wide FDs     sysctl kern.openfiles kern.maxfiles kern.maxfilesperproc
  socket buffers      netstat -m
                      sysctl kern.ipc.maxsockbuf net.inet.tcp.sendspace \
                             net.inet.tcp.recvspace kern.ipc.somaxconn
  I/O design check    procstat -k <pid>
\end{lstlisting}

\paragraph{OS tuning applied.}  FreeBSD's defaults ($\sim$64\,k sockets,
$\sim$29\,k FDs per proc, \texttt{SOMAXCONN}=128) cannot support 70\,k connections.
The tunables raised before the runs are summarised in Table~\ref{tab:bonus-tuning};
full rationale is in \texttt{bonus/README.md}.

\begin{table}[H]
\centering
\small
\begin{tabular}{@{}lll@{}}
\toprule
\textbf{Where} & \textbf{Setting} & \textbf{Reboot?} \\ \midrule
\texttt{/boot/loader.conf} & \texttt{kern.ipc.maxsockets=400000}      & Yes \\
\texttt{/boot/loader.conf} & \texttt{kern.ipc.nmbclusters=500000}     & Yes \\
\texttt{/etc/sysctl.conf}  & \texttt{kern.ipc.somaxconn=65535}        & No  \\
\texttt{/etc/sysctl.conf}  & \texttt{kern.maxfiles=300000}            & No  \\
\texttt{/etc/sysctl.conf}  & \texttt{kern.maxfilesperproc=250000}     & No  \\
\texttt{/etc/sysctl.conf}  & \texttt{net.inet.ip.portrange.first=10000} & No \\
\texttt{/etc/sysctl.conf}  & \texttt{net.inet.ip.portrange.last=65535}  & No \\
Per-shell (server + gen.)  & \texttt{ulimit -n 250000}                & No  \\
\texttt{lo0} aliases       & \texttt{ifconfig lo0 alias 127.0.0.\{2,3,4\}/32} & No \\
\bottomrule
\end{tabular}
\caption{OS-level tunables applied for the bonus experiment.  Applied once as root on
the FreeBSD~14 VM; see \texttt{bonus/README.md} \S3--\S5 for the exact commands and
persistence recipes.}
\label{tab:bonus-tuning}
\end{table}

%----------------------------------------------------------------------
\subsection{Changes made to the Exchange Server}
\label{sec:bonus-listen}
%----------------------------------------------------------------------

\textbf{Exactly one line} of the Exchange Server was changed for the bonus, and only
after the 40\,k run made the need explicit.  The original code called

\begin{lstlisting}
listen(server_fd, SOMAXCONN);   /* SOMAXCONN is the C-header macro = 128 on FreeBSD 14 */
\end{lstlisting}

On FreeBSD~14 the C header defines \texttt{SOMAXCONN} as \textbf{128},
\emph{independent} of the runtime sysctl \texttt{kern.ipc.somaxconn} that Table~\ref{tab:bonus-tuning}
raises to 65\,535.  Under the sustained $\sim$55\,000/s SYN flood produced by
\texttt{idle\_clients.py}, the 128-entry syncache overflowed during the 40\,k run and
\textbf{199 out of 40\,000} SYNs were RST'd by the kernel in a single burst starting at
attempt \#30\,994 (see Fig.~\ref{fig:bonus-40k-2}).  The fix is one character-count
change:

\begin{lstlisting}
if (listen(server_fd, 65535) == -1) { ... }   /* was: SOMAXCONN */
\end{lstlisting}

The kernel silently clamps the argument to \texttt{kern.ipc.somaxconn}, so
\texttt{65535} is safe on any FreeBSD.  \textbf{Scope of the change:} exactly one
syscall argument, sizing the syncache; the accept path, \texttt{kqueue} loop, message
framing and order-matching code are byte-for-byte identical.  The 10\,k / 20\,k / 30\,k
/ 40\,k rows below were collected with the original \texttt{SOMAXCONN} build; every row
from 50\,k onward uses the new \texttt{listen(65535)} build.  The direct A/B is
$40\text{k} \to 199$ lost vs.\ $50\text{k} \to 0$ lost -- the change eliminated the
losses even at higher load.  No other server-side change was made.

%----------------------------------------------------------------------
\subsection{Client-side source-IP rotation (for \texorpdfstring{$N > 55\,000$}{N > 55000})}
\label{sec:bonus-srcip}
%----------------------------------------------------------------------

A single \texttt{idle\_clients.py} process connecting from \texttt{127.0.0.1} to
\texttt{127.0.0.1:5000} cannot exceed about \textbf{55\,000 connections} in one run.
That is a TCP 4-tuple constraint, not a server limit: with $(src\_ip, src\_port,
dst\_ip, dst\_port)$ required to be unique and three of the four values fixed, only
$src\_port$ can vary, and only $\sim$55\,000 ephemeral ports exist even after widening
\texttt{net.inet.ip.portrange}.

To reach 60\,k and 70\,k in a single process, \texttt{idle\_clients.py} was extended
with three \emph{optional} flags -- \texttt{-{}-src-ips}, \texttt{-{}-src-cidr},
\texttt{-{}-src-count} -- that round-robin \texttt{s.bind((src\_ip, 0))} before
\texttt{connect()}.  Each distinct \texttt{src\_ip} has its own ephemeral-port space, so
$K$ source IPs raise the single-run ceiling to $\sim K \times 55\,000$.  Default
behaviour (all earlier runs, 10\,k--50\,k) is unchanged.

\paragraph{FreeBSD prerequisite.}  Unlike Linux, FreeBSD does not accept
\texttt{bind((127.x.x.x, 0))} unless the address is explicitly aliased on \texttt{lo0}
(\texttt{EADDRNOTAVAIL} otherwise).  Once, as root:

\begin{lstlisting}
ifconfig lo0 alias 127.0.0.2/32
ifconfig lo0 alias 127.0.0.3/32
ifconfig lo0 alias 127.0.0.4/32
\end{lstlisting}

The 60\,k and 70\,k rows below therefore invoke
\texttt{-{}-src-cidr 127.0.0.0/8 -{}-src-count 4}, drawing from
\{\texttt{127.0.0.1}, \texttt{127.0.0.2}, \texttt{127.0.0.3}, \texttt{127.0.0.4}\}.
The server side sees no difference from the earlier rows -- it is still just
\texttt{accept()}ing loopback connections.

%----------------------------------------------------------------------
\subsection{Full measurement table}
\label{sec:bonus-table}
%----------------------------------------------------------------------

All rows below come from \texttt{measure.sh} run against a live server that
\texttt{idle\_clients.py} is holding open at exactly $N$ established connections.
Column definitions match the Section~6.9 handout: server \texttt{RSS} and \%CPU from
\texttt{ps}, server FDs from \texttt{procstat -f}, system-wide FDs from
\texttt{sysctl kern.openfiles}, socket-buffer usage from \texttt{netstat -m} (mbuf
clusters in use / cap), and the established count cross-checked between
\texttt{sockstat} (minus one for the listening socket) and \texttt{idle\_clients.py}'s
own counter.

\begin{table}[H]
\centering
\scriptsize
\setlength{\tabcolsep}{4pt}
\begin{tabular}{@{}rrrrrrrr@{}}
\toprule
\textbf{$N$} & \textbf{RSS} & \textbf{\%CPU} & \textbf{Server FDs} &
\textbf{Sys FDs (used/max)} & \textbf{mbuf clusters (used/max)} &
\textbf{Net alloc} & \textbf{Established} \\
\midrule
10\,000 &  6\,140\,KB & 0.0\,\% & 10\,009 &  20\,094 / 300\,000 & 2\,049 / 500\,000 & 5\,644\,KB & 10\,000 / 10\,000 \\
20\,000 &  9\,996\,KB & 0.0\,\% & 20\,008 &  40\,092 / 300\,000 & 2\,051 / 500\,000 & 5\,648\,KB & 19\,999 / 20\,000$^{\dagger}$ \\
30\,000 & 13\,968\,KB & 0.0\,\% & 30\,009 &  60\,176 / 300\,000 & 2\,049 / 500\,000 & 5\,644\,KB & 30\,000 / 30\,000 \\
40\,000 & 17\,676\,KB & 0.0\,\% & 39\,806 &  79\,774 / 300\,000 & 2\,050 / 500\,000 & 5\,646\,KB & 39\,797 / 40\,000$^{\dagger}$ \\
50\,000 & 21\,888\,KB & 0.0\,\% & 50\,009 & 100\,197 / 300\,000 & 2\,049 / 500\,000 & 5\,644\,KB & 50\,000 / 50\,000$^{\ddagger}$ \\
60\,000 & 25\,696\,KB & 0.2\,\% & 60\,009 & 120\,198 / 300\,000 & 2\,050 / 500\,000 & 5\,646\,KB & 60\,000 / 60\,000$^{\ddagger\S}$ \\
\textbf{70\,000} & \textbf{29\,512\,KB} & \textbf{0.1\,\%} & \textbf{70\,009} & \textbf{140\,198 / 300\,000} & \textbf{2\,049 / 500\,000} & \textbf{5\,644\,KB} & \textbf{70\,000 / 70\,000}$^{\ddagger\S}$ \\
\bottomrule
\end{tabular}
\caption{Resource footprint of the Exchange Server holding $N$ idle TCP connections,
$N = 10\,000 \dots 70\,000$.  $^{\dagger}$ Original \texttt{listen(SOMAXCONN)} build;
transient listen-backlog RSTs during the SYN flood (1 at 20\,k, 199 at 40\,k) --
see \S\ref{sec:bonus-listen}.  $^{\ddagger}$ \texttt{listen(65535)} build.
$^{\S}$ Source-IP rotation across \texttt{127.0.0.\{1..4\}} (see \S\ref{sec:bonus-srcip}).
Column definitions in \S\ref{sec:bonus-table}.}
\label{tab:bonus-full}
\end{table}

\paragraph{Growth per 10\,k connections.}  Reading across
Table~\ref{tab:bonus-full}: server \texttt{RSS} grows by
$\approx3.9$\,MB per 10\,000 connections, i.e.\ $\sim$\textbf{390\,B/conn} of
user-space memory (essentially the size of \texttt{struct ClientConnection} +
one \texttt{conns[]} slot).  Server FDs grow by exactly one per connection.
System-wide FDs grow by two per connection because loopback holds both the
client end and the accepted server end.  \emph{mbuf clusters do not grow at
all}: they stay at $\sim$2\,050 across every row -- idle sockets \emph{reserve}
socket-buffer space but do not \emph{allocate} clusters until data flows.
\%CPU is 0 at steady state at every $N$; the server thread spends its entire
life parked inside a single \texttt{kevent()} call
(\texttt{procstat -k} confirms the kernel stack
\texttt{mi\_switch $\to$ sleepq\_wait\_sig $\to$ kqueue\_scan $\to$ sys\_kevent}
at every scale).

%----------------------------------------------------------------------
\subsection{Screenshots}
%----------------------------------------------------------------------

Three representative rows are shown as raw \texttt{measure.sh} output: the
smallest ($N=10\,000$), the row that exposed the listen-backlog bug
($N=40\,000$), and the target ($N=70\,000$).  The remaining rows are in
\texttt{bonus/} as \texttt{run\_20k.txt} through \texttt{run\_60k.txt}.

\shot{bonus/10k/1.jpeg}{$N=10\,000$: \texttt{measure.sh} output part~1 --
\texttt{ps}, \texttt{procstat -f}, \texttt{sysctl kern.openfiles}.
\texttt{RSS}=6\,140\,KB, \%CPU=0.0, server FDs = 10\,009,
system-wide FDs = 20\,094 / 300\,000.}{fig:bonus-10k-1}

\shot{bonus/10k/2.jpeg}{$N=10\,000$: \texttt{netstat -m} and socket-buffer
sysctls -- 2\,049 mbuf clusters in use / 500\,000 cap; 5\,644\,KB allocated to
network.  All 10\,000 idle connections consume $<0.5\,\%$ of the mbuf cluster
pool.}{fig:bonus-10k-2}

\shot{bonus/10k/3.jpeg}{$N=10\,000$: \texttt{procstat -k} of the server
thread and \texttt{sockstat} count.  The single thread is parked at
\texttt{sys\_kevent} (kqueue idle wait); \texttt{sockstat} sees 10\,001
sockets (10\,000 clients + 1 listener), confirming the connect count.}
{fig:bonus-10k-3}

\shot{bonus/40k/1.jpeg}{$N=40\,000$ (original \texttt{listen(SOMAXCONN)}
build): \texttt{ps} + \texttt{procstat -f}.
\texttt{RSS}=17\,676\,KB, server FDs = 39\,806, system-wide FDs = 79\,774 /
300\,000.}{fig:bonus-40k-1}

\shot{bonus/40k/2.jpeg}{$N=40\,000$: \texttt{idle\_clients.py} output during
the run -- 199 \texttt{ECONNRESET} failures in a single burst starting at
attempt \#30\,994.  Cause: syncache overflow against the C-header
\texttt{SOMAXCONN}=128, unrelated to the tuned sysctl (65\,535).  Fixed by
the one-line change in \S\ref{sec:bonus-listen}.}{fig:bonus-40k-2}

\shot{bonus/40k/3.jpeg}{$N=40\,000$: \texttt{netstat -m} -- 2\,050 mbuf
clusters in use, essentially unchanged from the 10\,k row.  Idle
connections still do not consume clusters even after $4\times$ scale-up.}
{fig:bonus-40k-3}

\shot{bonus/40k/4.jpeg}{$N=40\,000$: \texttt{procstat -k} + \texttt{sockstat}.
Kernel stack is still parked at \texttt{sys\_kevent} despite 39\,797
established connections; \%CPU still 0.0.}{fig:bonus-40k-4}

\shot{bonus/70k/1.jpeg}{$N=70\,000$ (target): \texttt{ps} + \texttt{procstat -f}.
\texttt{RSS}=29\,512\,KB (only $\sim$29\,MB for 70\,000 idle connections),
server FDs = 70\,009, system-wide FDs = 140\,198 / 300\,000 (47\,\% of
cap).}{fig:bonus-70k-1}

\shot{bonus/70k/2.jpeg}{$N=70\,000$: \texttt{idle\_clients.py} final
progress -- \textbf{70\,000 established / 0 failed} with source-IP rotation
across \texttt{127.0.0.\{1..4\}} and the \texttt{listen(65535)} build.}
{fig:bonus-70k-2}

\shot{bonus/70k/3.jpeg}{$N=70\,000$: \texttt{netstat -m} -- 2\,049 mbuf
clusters in use, \emph{identical} to the 10\,k row.  Confirms that idle
sockets consume zero mbuf clusters at any scale.}{fig:bonus-70k-3}

\shot{bonus/70k/4.jpeg}{$N=70\,000$: \texttt{procstat -k} + \texttt{sockstat}
= 70\,001 (70\,000 clients + 1 listener).  Server thread still parked at
\texttt{sys\_kevent}; \%CPU 0.1 (tail of the connect flood, not
steady-state).}{fig:bonus-70k-4}

%----------------------------------------------------------------------
\subsection{Answers to Section~6.9 questions}
%----------------------------------------------------------------------

\qbox{Q1.\ At each connection count, is the server able to maintain all requested
connections?  If not, at which count does it first fail?}

\begin{answer}
Yes at every measured count. The Exchange Server successfully maintained
10\,000 / 20\,000 / 30\,000 / 40\,000 / 50\,000 / 60\,000 / 70\,000 simultaneous
idle TCP connections; the established count equals or exceeds the request in
every row of Table~\ref{tab:bonus-full}.  Two transient RST bursts were
observed \emph{during the connect flood} (1 loss at 20\,k, 199 losses at
40\,k), both caused by the FreeBSD C header \texttt{SOMAXCONN}=128 capping the
listen syncache below what the tuned \texttt{kern.ipc.somaxconn}=65\,535 sysctl
allowed.  After the one-line fix (\S\ref{sec:bonus-listen}) every higher-$N$
run (50\,k / 60\,k / 70\,k) completed with zero failures.  No ``first-fails''
count was reached; the experiment stopped at the assignment's target of
70\,000 with the server still idle at 0\,\% CPU and $\sim$47\,\% of the
system-wide FD cap unused.
\end{answer}

\qbox{Q2.\ What is the first significant bottleneck?}

\begin{answer}
\textbf{Kernel socket-structure and socket-buffer-cluster pools, not anything
under application control.}  With FreeBSD defaults the first two ceilings hit
were:
\begin{enumerate}[leftmargin=1.5em,itemsep=0.15em]
  \item \texttt{kern.ipc.maxsockets} ($\approx$64\,177 on this VM). Loopback
        consumes two kernel sockets per connection, so the effective ceiling
        is $\sim$32\,000 -- exactly where \texttt{ECONNRESET} first appeared
        in an untuned exploratory run.
  \item \texttt{kern.ipc.nmbclusters} ($\approx$125\,007). The next failure
        mode would have been \texttt{ENOBUFS} on \texttt{socket()} itself.
\end{enumerate}
Both are boot-time tunables raised via \texttt{/boot/loader.conf}
(Table~\ref{tab:bonus-tuning}).  Runtime-tunable ceilings on
\texttt{kern.maxfiles}, \texttt{kern.maxfilesperproc},
\texttt{kern.ipc.somaxconn} and \texttt{net.inet.ip.portrange.*} were also
raised.  Once these were lifted, \emph{the server itself} never became the
bottleneck in the 10\,k--70\,k range: \texttt{RSS} stayed under 30\,MB,
\%CPU at 0, and \texttt{mbuf clusters in use} flat at $\sim$2\,050 regardless
of $N$ (Table~\ref{tab:bonus-full}).  The next-ceiling ranking, if scaled
further without more tuning, is:
\texttt{maxsockets} $\to$ \texttt{nmbclusters} $\to$ FD limits
$\to$ ephemeral ports on the \emph{generator} side ($\sim$55\,k per source
IP, addressed by source-IP rotation, \S\ref{sec:bonus-srcip})
$\to$ reserved socket-buffer memory
(\texttt{sendspace}+\texttt{recvspace} = 96\,KB $\times$ 2 sockets/conn).
\end{answer}

\qbox{Q3.\ How does the concurrency / I/O design contribute to the observed
bottleneck?}

\begin{answer}
\textbf{It contributes essentially zero.}  The Exchange Server uses a
single-threaded event loop over \texttt{kqueue} with non-blocking sockets --
the correct primitive on FreeBSD for high-fanout servers.  Concretely:
\begin{itemize}[leftmargin=1.5em,itemsep=0.15em]
  \item \emph{One thread, always parked in \texttt{kevent()}.}
        \texttt{procstat -k} shows the kernel stack
        \texttt{mi\_switch $\to$ sleepq\_wait\_sig $\to$ kqueue\_scan $\to$
        sys\_kevent} at every $N$ -- no per-connection thread stack, no
        \texttt{fd\_set} scan (which would cap us at \texttt{FD\_SETSIZE}),
        no busy-poll; \%CPU is 0 when no events fire.
  \item \emph{Per-connection application state is $\sim$360\,B}
        (\texttt{struct ClientConnection}) plus one pointer slot in
        \texttt{conns[]}; the write buffer is \texttt{NULL} until data is
        actually queued.  \texttt{RSS} scaled at $\sim$390\,B/conn across
        the whole range -- dominated by that struct.
  \item \emph{Connection table doubles.}  \texttt{conns[]} grows via
        \texttt{realloc} from \texttt{MIN\_CAPACITY}=10; at $N=70\,000$ it
        occupies only $\sim$640\,KB.
\end{itemize}
Because the design has no $O(N)$ work per event and no per-connection
resource beyond that small struct, none of the bottlenecks in Q2 can be
attributed to it -- they are all kernel-side per-connection allocations
that any correct TCP server pays.
\end{answer}

\qbox{Q4.\ Would changing the I/O / concurrency mechanism help with the
bottleneck identified?}

\begin{answer}
\textbf{No -- the current design is already the least-costly alternative.}
The Q2 bottlenecks (\texttt{maxsockets}, \texttt{nmbclusters},
\texttt{maxfiles}) are kernel-side per-connection allocations paid by any
correct TCP server regardless of user-space multiplexing.  Concretely:
\begin{itemize}[leftmargin=1.5em,itemsep=0.15em]
  \item \emph{Thread-per-connection} would collapse first: each pthread
        costs a $\sim$1--2\,MB stack, a \texttt{struct thread}, and a slot
        in \texttt{kern.threads.max\_threads\_per\_proc} ($\sim$1\,500 by
        default).  At 70\,k connections it would exhaust threads long
        before touching the kernel-socket caps we hit.
  \item \emph{\texttt{select()}} caps at \texttt{FD\_SETSIZE}=1024 and
        performs $O(N)$ scans per event; worse than \texttt{kqueue} even
        at 10\,k.
  \item \emph{\texttt{poll()}} avoids the \texttt{FD\_SETSIZE} cap but still
        scans $O(N)$ on every event; at 70\,k it would use substantially
        more CPU than \texttt{kqueue}, though it could still scale in
        principle.
\end{itemize}
Any alternative therefore either (a) hits the same kernel-resource ceiling
while consuming more CPU/memory, or (b) fails earlier for its own reasons.
\texttt{kqueue} is the right tool; the ceiling we reached is below what
\texttt{kqueue} is capable of.
\end{answer}

\qbox{Q5.\ Quantify the trade-off -- if you change your implementation or
configuration, measure the relevant resource again and compare.}

\begin{answer}
Exactly one configuration parameter was changed (the \texttt{listen()}
backlog, \S\ref{sec:bonus-listen}) with all other tuning and the client
generator held constant.  The relevant resource is
``successful vs.\ failed connects during the flood'', measured at $N=40\,000$
before the change and $N=50\,000$ after (higher load):

\begin{center}
\small
\begin{tabular}{@{}llrr@{}}
\toprule
\textbf{Regime} & \textbf{Setting} & \textbf{Established / requested} & \textbf{Lost} \\ \midrule
Before & \texttt{listen(server\_fd, SOMAXCONN)}, \texttt{SOMAXCONN}=128 & 39\,797 / 40\,000 & 199 (0.50\,\%) \\
After  & \texttt{listen(server\_fd, 65535)} & 50\,000 / 50\,000 & 0 (0.00\,\%) \\
\bottomrule
\end{tabular}
\end{center}

Higher load, same tuning, same generator settings, same measurement recipe --
the one-line change eliminated the losses entirely, and 60\,k and 70\,k
subsequently completed with 0 failures as well (Table~\ref{tab:bonus-full}).
\textbf{Runtime cost of the change: zero.}  Raising the backlog does not
allocate anything eagerly; it only sets the cap on how many half-open
connections the syncache is allowed to hold.  No change to accepted-connection
memory, CPU, latency, or protocol behaviour.
\end{answer}

%----------------------------------------------------------------------
\subsection{Notes and caveats}
%----------------------------------------------------------------------

\begin{itemize}[leftmargin=1.5em,itemsep=0.2em]
  \item \textbf{Loopback doubles the kernel-socket cost.} Every connection
        in this experiment holds two kernel sockets (client end + accepted
        server end) because both live in the same VM.  A real deployment
        driven from a second host would need roughly half the
        \texttt{kern.ipc.maxsockets} headroom for the same $N$.
  \item \textbf{The 40\,k row is intentionally left in the table with its
        199 losses.}  It is the direct evidence for the
        \texttt{SOMAXCONN} bug and the A/B pair for Q5.  Re-running it on
        the fixed build would also yield 0 losses (as 50\,k / 60\,k / 70\,k
        already show).
  \item \textbf{Idle vs.\ active is a very different workload.}  All rows
        above measure \emph{idle} connections, per the handout.  Active
        traffic would allocate mbuf clusters proportional to
        in-flight bytes; \texttt{netstat -m} was flat here only because no
        data flowed.  The order-book code path itself was exercised in
        Experiments~2--8; scaling that under load is out of scope for
        Section~6.9.
  \item \textbf{Non-zero \%CPU at 60\,k and 70\,k} (0.2\,\% and 0.1\,\%) is
        the tail of the connect flood being sampled before it fully
        subsided, not a steady-state cost of holding the connections.
        \texttt{procstat -k} confirmed the server was already back in
        \texttt{sys\_kevent} at snapshot time.
  \item \textbf{Reproducibility.}  Every command and setting used to
        produce Table~\ref{tab:bonus-full} is listed in
        \texttt{compiled/bonus/README.md}; per-run raw outputs are stored
        as \texttt{compiled/bonus/\{10k,20k,\dots,70k\}/*.txt} in the
        submission alongside the screenshots reproduced in this appendix.
\end{itemize}

\end{document}
